\specialsection{Введение}

В настоящее время количество доступной информации постоянно увеличивается.
Современные распределённые системы хранения должны одновременно обеспечивать
высокую доступность данных, устойчивость к отказам оборудования и эффективное
использование дискового пространства. Эти требования особенно важны для
крупных хранилищ, где данные размещаются на множестве серверов, а отказы
отдельных дисков или узлов являются не исключением, а штатной частью
эксплуатации \cite{shen2024erasurecoding}. В таких условиях способ
организации избыточности напрямую влияет на стоимость хранения, задержку
чтения и цену восстановления после отказов.

На практике широко используются два основных подхода к отказоустойчивому
хранению: репликация и коды, исправляющие стирания. Репликация проста в
реализации и обеспечивает дешёвое чтение при отказах, но требует большого
расхода дискового пространства. Коды с исправлением стираний, напротив,
позволяют существенно снизить объём избыточных данных, но делают чтение и
восстановление при отказах более дорогими
\cite{huang2012azure,muralidhar2014f4}.

На практике данные нередко сначала хранятся с репликацией, а затем переводятся
в схемы с кодами стираний для снижения расхода места
\cite{li2017ear,kim2024morph}. Однако такой переход сам по себе может быть
дорогим: в распространённой реализации перекодирование сводится к чтению
данных, записи нового представления и удалению старого \cite{kim2024morph}.
Сложность выбора усиливается тем, что нагрузка на данные обычно неравномерна:
разные объекты и участки объектов имеют разную интенсивность чтения, которая
изменяется со временем \cite{levandoski2013hotcold,muralidhar2014f4}.

В работах о хранилищах такую активность часто описывают через температуру
данных: горячие данные получают много запросов, а более холодные данные
читаются реже \cite{muralidhar2014f4,song2024hsm}. В данной работе температура
используется как численная метрика текущей интенсивности чтений, а не только
как принадлежность объекта к одному из заранее заданных классов. При высокой
активности важнее низкая ожидаемая стоимость чтения, при низкой ---
компактность хранения. Кроме того, меняется состояние самого кластера: при
достаточном свободном месте система может предпочитать более быстрые схемы, а
при росте заполненности дисков --- более компактные. Следовательно,
фиксированная схема избыточности может быть либо слишком дорогой по месту, либо
слишком дорогой по чтению, восстановлению и последующему перекодированию.

Таким образом, актуальной становится задача не только выбора одной схемы
избыточности при записи данных, но и управления схемой хранения в течение
жизненного цикла данных. Такая задача требует одновременно учитывать текущую
активность чтений, доступное дисковое пространство и стоимость фоновых
переходов, поскольку само перекодирование может конкурировать с
пользовательской нагрузкой за дисковые и сетевые ресурсы.

В данной работе рассматривается подход, в котором схема хранения является
изменяемым состоянием жизненного цикла данных. Основная идея состоит в том,
чтобы описывать разные формы избыточности в едином пространстве состояний и
выбирать фоновые переходы между ними по функции полезности. В отличие от
фиксированной цепочки переходов, планировщик учитывает текущую температуру
данных, занятое место, ожидаемую стоимость чтения и стоимость фоновой работы.
Формальная модель такого пространства состояний вводится в первой главе.

Работа ориентирована на append-нагрузку: запись новых объектов и дозапись
существующих без изменения уже закрытых полос данных.

Цель работы --- разработать модель и архитектурный подход адаптивного
отказоустойчивого хранения данных с дозаписью, в котором система может
автоматически менять форму избыточности между состояниями с преобладанием
реплик и состояниями с преобладанием проверочных блоков в зависимости от
активности данных, занятого места, ожидаемой стоимости чтения и стоимости
фонового перехода между схемами.

Для достижения цели были поставлены следующие задачи:
\begin{enumerate}
    \item рассмотреть существующие подходы к отказоустойчивому хранению:
    репликацию, коды стираний, локально восстанавливаемые и конвертируемые
    коды;
    \item рассмотреть существующие подходы к учёту активности данных и
    моделированию температуры как численной метрики активности чтений;
    \item сформулировать модель гибридной схемы хранения, объединяющей
    реплики, локальные и глобальные проверочные блоки;
    \item определить допустимые переходы между схемами, сохраняющие заданный
    уровень отказоустойчивости;
    \item разработать модель оценки полезности перехода с учётом занятого
    места, ожидаемой нагрузки чтения и стоимости фонового перекодирования;
    \item спроектировать архитектуру фонового перекодировщика и планировщиков,
    выбирающих переходы между схемами;
    \item реализовать экспериментальную систему, поддерживающую запись
    объектов, чтение при отказах и фоновые переходы между схемами;
    \item провести экспериментальную оценку поведения системы в сценариях
    охлаждения данных, неравномерной нагрузки чтения и отказов.
\end{enumerate}

Дальнейшая структура работы организована следующим образом. В разделе
<<Обзор литературы>> рассматриваются существующие подходы к отказоустойчивому
хранению и работы, связанные с изменением схем хранения. В первой главе
вводятся теоретические основы модели: репликация, коды Рида--Соломона
\cite{reed1960polynomial}, локально восстанавливаемые коды
\cite{gopalan2012lrc,huang2012azure}, гибридные схемы, сочетающие репликацию
и коды стираний \cite{lucani2025hyres}, температура как метрика интенсивности
чтений, модель ожидаемой стоимости чтения и функция полезности перехода. Во
второй главе описывается архитектура решения: размещение блоков, фоновые
переходы между схемами, чтение при отказах и работа планировщиков. В третьей
главе рассматриваются детали реализации системы. В четвёртой главе приводятся
результаты экспериментов и сравнение с фиксированными подходами. В заключении
подводятся итоги работы, фиксируются полученные результаты, ограничения решения
и направления дальнейшего развития.

\pagebreak
